The goal of this paper is to compare Operational Transformation and CRDTs as approaches for handling
conflicts in distributed systems. Since both approaches are well documented in research literature and in
the documentation of existing software, this comparison is done by looking at existing sources rather than
running our own tests or experiments. We look at what each approach promises, how it works in theory, and
how real systems have implemented it.

\subsection{Comparison Criteria}

To make the comparison structured and fair, we evaluate both approaches based on the same set of criteria:

\begin{enumerate}
    \item \textbf{Response time:} Can a node apply a change immediately, or does it have to wait for a
    response from another server first? Fast response is important for a smooth user experience.
    \item \textbf{Consistency:} Do all nodes end up with the same data after syncing? This is the most
    important requirement for any conflict resolution system \cite{shapiro2011comprehensive}.
    \item \textbf{Difficulty of implementation:} How hard is it to build the approach correctly? OT in
    particular has a history of subtle bugs that are easy to introduce \cite{nichols1995jupiter}.
    \item \textbf{Suitability for decentralised use:} Does the approach require a central server, or can
    it work in a peer-to-peer setup where nodes are sometimes offline?
    \item \textbf{Storage and network cost:} How much extra data needs to be stored or sent over the
    network to make the conflict resolution work?
\end{enumerate}

\subsection{Test Scenarios}

We apply these criteria to two typical situations that often cause problems in distributed systems:

\textbf{Many changes at once:} Several nodes make changes to the same data at the same time, without
knowing about each other's changes yet. This is the core problem that both OT and CRDTs are trying to solve.

\textbf{Offline and reconnect:} A node is temporarily disconnected from the network, continues to accept
changes locally, and then needs to sync everything when it comes back online. This is a very common
real-world situation -- for example when editing a document on a phone without internet access
\cite{automerge2024, nicolaescu2015}.

\subsection{Systems Used as Examples}

To show how the theory works in practice, we look at a few real systems that use either OT or CRDTs:

\begin{itemize}
    \item \textbf{Jupiter} \cite{nichols1995jupiter} -- an early collaborative editing system that used OT
    with a central server.
    \item \textbf{Riak KV} \cite{riak2024} -- a database that uses CRDT data types to handle conflicting
    writes automatically.
    \item \textbf{Redis Active-Active} \cite{redis2024} -- a database system that allows writes in multiple
    locations at the same time, and uses CRDTs to merge them.
    \item \textbf{Automerge} \cite{automerge2024} and \textbf{Yjs} \cite{nicolaescu2015} -- two
    open-source libraries that use CRDTs to let developers build collaborative, offline-capable
    applications.
\end{itemize}

For each system we only look at the official documentation and published descriptions -- no new tests
were run.
